\documentclass[
%	draftmark = true,   % 草稿模式下, 每页底部将打印相关版本信息.
%	traditional = true,
	colors = true,
	coverpage = coverpage.tex,
%	coverpage = coverpage.pdf,
    geometry = b5,
%	index = true
]{PRbook}

\usepackage{mathrsfs}
\usepackage{stmaryrd} \SetSymbolFont{stmry}{bold}{U}{stmry}{m}{n}	% 避免警告 (stmryd 不含粗体故)
\usepackage{array}
\usepackage{makecell}	% 便于制表
\usepackage{tikz-cd}  % 使用 TikZ 绘图
\usetikzlibrary{positioning, patterns, calc, matrix, shapes.arrows, shapes.symbols, fit, arrows.meta, braids}
\usepackage{pgfplots}
\usepackage{graphicx}
\pgfplotsset{compat=newest}

\usepackage{myarrows}				% 使用自定义的可伸缩箭头
\usepackage{mycommand}				% 引入自定义的惯用的命令



% 公式和图片按节编号
\numberwithin{equation}{section}
\numberwithin{figure}{section}

\usepackage[unicode, bookmarksnumbered]{hyperref}	% 启动超链接和 PDF 文档信息所需

\hypersetup{
	pdfauthor = {张竣淞(Zhang Junsong)},
	%pdftitle = {****},
	%pdfkeywords = {****},
	CJKbookmarks = true}

\usepackage[open, openlevel=0]{bookmark}	% 优化 PDF 书签管理

\addbibresource{ref.bib}

\begin{document}

\frontmatter	% 制作封面和目录.

\mainmatter		% 正文开始：逐章引入 TeX 代码
这是一份\LaTeX 模板，是基于李文威老师的著作《代数学方法：卷二》的源代码进行简化修改的，主要是修改了一些字体的配置，
宋体改为使用思源宋体（Source Han Serif SC），黑体使用思源黑体（Source Han Sans SC）,由于可变字体会出现笔者无法
修复的问题，所以均使用静态字体，这两个系列字体均为开源字体可自行下载，为避免出现问题，建议选择“为全部用户安装”；仿宋和楷体则改为使用Windows系统自带的仿宋和楷体。\par
修改后的配置在VS Code中使用XeLatex+Biber+XeLatex+XeLatex编译链.\par
另外，李文威老师的自定义箭头和命令主要应用于代数领域，个性较强，笔者可能使用较少，但仍然保留以备不时之需。\par
最后在此表示对李文威老师的钦佩与感谢，由于笔者对开源协议不太熟悉，如有不妥之处，烦请及时告知，以免损害李文威老师或其它人的合法权益。
\part{示例}
\chapter{环境示例}
这一章主要展示自定义的定理、定义等\emph{环境}以及编号规则。
\begin{tips}
环境种类较多，如有遗漏，可参见cls文件。
\end{tips}
\section{正文环境}
首先展示正文中涉及到的例子。\par
这是一个术语\emph{索引}的展示，关键字为“索引”。\par
\index{soyin@索引（index）}
这是一个符号索引的展示，关键字为“$\mathbb{R}$”。\par
\index[sym]{Prob@$\mathbb{R}$}
这是一个引用的例子，\cite{AM10}.
\begin{theorem}
    Brown Motion存在。
\end{theorem}

\begin{proof}
    此处略。
\end{proof}




\bookmarksetup{startatroot}	% PDF 书签复位, 不再作为 \part 的子节点.
%\appendix
%\include{appendix1}
%\include{appendix2}

\backmatter
% 使用 bibLaTeX 制作书目
\printbibliography[heading=bibintoc]

\end{document}